% ============================================================
%  PART IV  --  My Data
% ============================================================
\part{My Data}

% ── Chapter 12: My Data Overview ──────────────────────────────
\chapter{My Data --- Overview}
\label{ch:mydata}

My Data is FairWindSK's built-in resource manager.
It provides a dedicated, touch-optimised interface for all navigational data stored on the Signal~K
server: waypoints, routes, regions, notes, charts, track history, and files.

Open My Data by tapping the \key{My Data} icon (≡) in the Bottom Bar, or by tapping its tile on the
launcher.

\screenshot{mydata-main-tabs}{My Data --- tab bar showing all seven tabs}{}

\begin{notebox}
All resources in My Data are stored on the Signal~K server, not on the FairWindSK device.
They are accessible from any device connected to the same Signal~K server simultaneously.
\end{notebox}

\section{My Data Tab Overview}

My Data is organised into seven tabs.
Each tab loads lazily --- it is created only the first time it is tapped, keeping startup fast.

\rowcolors{2}{LtGray}{white}
\begin{longtable}{C{0.7cm} L{2.8cm} L{10.2cm}}
\toprule
\rowcolor{Navy}
{\sffamily\bfseries\color{white} Tab} &
{\sffamily\bfseries\color{white} Name} &
{\sffamily\bfseries\color{white} Contents} \\
\midrule
0 & Waypoints & Named positions: anchorages, hazards, marinas, fishing marks, race marks. \\
1 & Routes & Ordered sequences of waypoints forming planned passages. \\
2 & Regions & Geographic boundary areas --- SAR search regions, exclusion zones, anchorage limits. \\
3 & Notes & Free-text annotations attached to positions. \\
4 & Charts & Chart tile source catalogue for installed chart providers. \\
5 & Tracks & Historical GPS track --- the vessel's recorded position history. \\
6 & Files & Direct filesystem browser for files accessible through Signal~K. \\
\bottomrule
\end{longtable}

\section{The Common Toolbar}

Tabs 0--4 (Waypoints, Routes, Regions, Notes, Charts) share a common interface layout:

\begin{itemize}
  \item A \textbf{title label} showing the resource type.
  \item A \textbf{search field} for real-time filtering of the list. Searching is case-insensitive and checks all visible columns simultaneously.
  \item A \textbf{scrollable table} showing all resources from the Signal~K server. Rows are 64\,px tall for comfortable touch selection. Alternating row colours aid readability. Column headers are clickable for sorting.
  \item A \textbf{toolbar} of icon buttons:
\end{itemize}

\rowcolors{2}{LtGray}{white}
\begin{longtable}{L{3cm} L{10.7cm}}
\toprule
\rowcolor{Navy}
{\sffamily\bfseries\color{white} Button} &
{\sffamily\bfseries\color{white} Action} \\
\midrule
Open ($\rightarrow$) & Opens the selected resource --- shows its details, opens an embedded map preview (Waypoints), or views the attachment (Notes). Accent-coloured; primary action. \\
Edit (✏) & Opens the resource editor drawer for the selected resource. \\
New (+) & Opens the resource editor drawer to create a new resource. \\
Delete (🗑) & Deletes the selected resource after confirmation. Cannot be undone. \\
Import & Imports resources from a GeoJSON or JSON file. File picker opens in the Bottom Bar drawer. \\
Export & Exports the selected resource (or all resources if none selected) to a GeoJSON file. \\
Refresh & Forces a reload of all resources from the Signal~K server. \\
\bottomrule
\end{longtable}

\begin{notebox}
On Raspberry Pi OS (ARM Linux), the toolbar uses text-only buttons to avoid icon rendering
issues on that platform.
The button layout and functionality are identical; only the visual presentation differs.
\end{notebox}

% ── Chapter 13: Waypoints ─────────────────────────────────────
\chapter{Waypoints}
\label{ch:waypoints}

Waypoints are named positions stored on the Signal~K server.
They are used for anchorages, hazards, marinas, fuel stops, fishing marks, race marks, and any
other position you want to record and navigate to.

\screenshot{mydata-waypoints-list}{My Data --- Waypoints tab, list view with name, description, type, lat/lon, timestamp}{}

\section{Waypoints Table Columns}

The waypoints table shows six columns:

\rowcolors{2}{LtGray}{white}
\begin{longtable}{L{2.4cm} L{11.3cm}}
\toprule
\rowcolor{Navy}
{\sffamily\bfseries\color{white} Column} &
{\sffamily\bfseries\color{white} Description} \\
\midrule
Name & Display name of the waypoint, taken from the \code{name} field or from \code{feature.properties.name} in the GeoJSON resource. \\
Description & Free-text description of the waypoint. \\
Type & The waypoint type tag, e.g.\ \code{alarm-mob} (POB waypoints created by FairWindSK), \code{anchor}, \code{marina}, or any custom type. POB waypoints are visually distinguishable by their type. \\
Latitude & Latitude in your configured coordinate format (right-aligned). \\
Longitude & Longitude in your configured coordinate format (right-aligned). \\
Timestamp & ISO 8601 creation or last-update timestamp. \\
\bottomrule
\end{longtable}

\section{Searching Waypoints}

The search field above the table filters waypoints by any visible column.
Type part of a name, description, type, or coordinate value.
The table updates in real time as you type.

\section{Creating a Waypoint}

Tap \key{New} to open the Waypoint Editor drawer.


\subsection{Editor Fields}

\rowcolors{2}{LtGray}{white}
\begin{longtable}{L{3cm} L{10.7cm}}
\toprule
\rowcolor{Navy}
{\sffamily\bfseries\color{white} Field} &
{\sffamily\bfseries\color{white} Description} \\
\midrule
Name & The display name shown in the table, Top Bar waypoint widget, and chart plotter. \\
Description & Free-text notes about this position. \\
Type & Optional type tag. Leave blank for a generic waypoint. Use \code{anchor} for anchoring positions, \code{hazard} for obstructions, etc. \\
Latitude & Degrees with decimal minutes or decimal degrees, depending on the active coordinate format. Range: $-90$ to $+90$. \\
Longitude & Range: $-180$ to $+180$. \\
Altitude & Altitude in metres (optional). Leave as 0 if not relevant. \\
\bottomrule
\end{longtable}

After filling in the fields, tap \key{OK} in the editor.
The waypoint is created immediately on the Signal~K server and appears in the table.

\section{Editing a Waypoint}

Select a waypoint in the table, then tap \key{Edit}.
The same editor drawer opens with the existing values pre-filled.
Modify any field and tap \key{OK}.
The change is written immediately to Signal~K.

\section{Deleting a Waypoint}

Select a waypoint and tap \key{Delete}.
A confirmation drawer appears.
Tap \key{Delete} to confirm.
The waypoint is removed from the Signal~K server permanently.

\begin{warningbox}
Deletion cannot be undone.
POB waypoints (type \code{alarm-mob}) should generally be retained for accident reporting and debrief,
not deleted.
\end{warningbox}

\section{Opening a Waypoint --- the Detail View}

Select a waypoint and tap \key{Open} (or double-tap the row).
The Waypoint Detail view opens.


The detail view shows:

\begin{itemize}
  \item Full name, description, type, and formatted coordinates.
  \item An embedded Freeboard-SK map preview centred on the waypoint's position, if Freeboard-SK is installed on the connected Signal~K server. The map uses a JavaScript function to compute the view bounding box from the waypoint coordinates and zoom level.
  \item A \key{Navigate} button to set the waypoint as the active navigation destination.
\end{itemize}

\begin{notebox}
The embedded Freeboard-SK preview is disabled on Raspberry Pi OS (ARM Linux) builds to avoid
performance issues with simultaneous WebEngine instances on that platform.
Coordinates and descriptions are always shown in text form regardless.
\end{notebox}

\section{Navigating to a Waypoint}

From the detail view, tap \key{Navigate}.
FairWindSK calls \path{navigateToWaypoint} on the Signal~K client, which writes the waypoint's
resource path (e.g.\ \path{/resources/waypoints/\{uuid\}}) as the active navigation destination on
Signal~K.
Immediately:

\begin{itemize}
  \item The Top Bar BTW, DTW, TTG, and ETA widgets begin showing values relative to this waypoint.
  \item Freeboard-SK and any other navigation application subscribed to \skpath{navigation.course} show the destination.
  \item The POB bearing and distance displays (when the POB bar is active) use this same destination.
\end{itemize}

\section{Importing Waypoints}

Tap \key{Import}. A file picker opens in the Bottom Bar drawer.
Select a GeoJSON file containing \code{Feature} objects with \code{Point} geometry.

FairWindSK:
\begin{enumerate}
  \item Reads the GeoJSON file.
  \item Extracts each \code{Point} feature as a waypoint.
  \item Creates or updates each waypoint on the Signal~K server (update if a matching UUID is found; create if not).
  \item Reports the count of imported waypoints in an informational drawer.
  \item Selects the first imported waypoint in the table.
\end{enumerate}

The import file filter accepts \code{.geojson} and \code{.json} files.

\section{Exporting Waypoints}

Tap \key{Export}. If a waypoint is selected, only that waypoint is exported.
If nothing is selected, all waypoints are exported.
A file-save picker opens. Choose a filename and location.
The output is a valid GeoJSON \code{FeatureCollection} with \code{Point} geometry features.

% ── Chapter 14: Routes ────────────────────────────────────────
\chapter{Routes}
\label{ch:routes}

Routes are ordered sequences of waypoints forming planned passages.
They are typically created in a chart plotter (Freeboard-SK) and stored on the Signal~K server.

\screenshot{mydata-routes-list}{My Data --- Routes tab, list view with name, geometry type, point count, timestamp}{}

\section{Routes Table Columns}

\rowcolors{2}{LtGray}{white}
\begin{longtable}{L{2.4cm} L{11.3cm}}
\toprule
\rowcolor{Navy}
{\sffamily\bfseries\color{white} Column} &
{\sffamily\bfseries\color{white} Description} \\
\midrule
Name & Display name of the route. \\
Description & Free-text description or summary. \\
Geometry & The GeoJSON geometry type --- typically \code{LineString} for simple routes, \code{MultiLineString} for multi-leg routes with gaps, or \code{sar-spiral} / \code{sar-parallel} for SAR patterns created by the POB feature. \\
Points & The number of coordinate points in the route geometry. Displayed centred. \\
Timestamp & ISO 8601 creation or last-update timestamp. \\
\bottomrule
\end{longtable}

\section{Opening a Route --- the Detail View}

Select a route and tap \key{Open}.


The detail view shows:
\begin{itemize}
  \item Name, description, geometry type, and point count.
  \item A GeoJSON geometry preview if available.
  \item An \key{Activate} button to set the route's first waypoint as the active navigation destination.
\end{itemize}

\section{Creating and Editing Routes}

\subsection{Creating a Route}

Tap \key{New}. The Route Editor drawer opens.

\rowcolors{2}{LtGray}{white}
\begin{longtable}{L{3cm} L{10.7cm}}
\toprule
\rowcolor{Navy}
{\sffamily\bfseries\color{white} Field} &
{\sffamily\bfseries\color{white} Description} \\
\midrule
Name & Display name of the route. \\
Description & Free-text notes. \\
Coordinates (JSON) & The GeoJSON coordinate array for the route geometry. Enter as a JSON array of \code{[longitude, latitude]} pairs. \\
Geometry (JSON) & Advanced: edit the full GeoJSON geometry object if needed. \\
Properties (JSON) & Advanced: additional GeoJSON feature properties. \\
\bottomrule
\end{longtable}

\subsection{Activating a Route}

From the route detail view, tap \key{Activate}.
FairWindSK sets the first waypoint of the route as the active Signal~K navigation destination.
The Autopilot Route mode (Chapter~\ref{ch:autopilot}) will then follow the active route automatically.

\section{Importing and Exporting Routes}

Import and export use the same GeoJSON \code{FeatureCollection} format as Waypoints.
The geometry type in the GeoJSON determines how the route is stored on Signal~K.

\begin{tipbox}
SAR search routes created by the POB feature (Chapter~\ref{ch:pob}) appear in the Routes tab
automatically after creation.
They have geometry types \code{sar-spiral} and \code{sar-parallel} and can be exported as GeoJSON
for voyage records.
\end{tipbox}

% ── Chapter 15: Regions ───────────────────────────────────────
\chapter{Regions}
\label{ch:regions}

Regions are geographic boundary areas defined as GeoJSON polygons.
They are used for SAR search areas (automatically created by the POB feature), marine protected area
boundaries, exclusion zones, anchorage limits, and fishing grounds.

\screenshot{mydata-regions-list}{My Data --- Regions tab, list view with name, geometry type, timestamp}{}

\section{Regions Table Columns}

\rowcolors{2}{LtGray}{white}
\begin{longtable}{L{2.4cm} L{11.3cm}}
\toprule
\rowcolor{Navy}
{\sffamily\bfseries\color{white} Column} &
{\sffamily\bfseries\color{white} Description} \\
\midrule
Name & Display name of the region. \\
Description & Free-text description. \\
Geometry & The GeoJSON geometry type --- typically \code{Polygon} or \code{MultiPolygon}. SAR regions use \code{Polygon}. \\
Timestamp & ISO 8601 creation or last-update timestamp. \\
\bottomrule
\end{longtable}

\section{Creating a Region}

Tap \key{New}. The Region Editor drawer opens.

\rowcolors{2}{LtGray}{white}
\begin{longtable}{L{3cm} L{10.7cm}}
\toprule
\rowcolor{Navy}
{\sffamily\bfseries\color{white} Field} &
{\sffamily\bfseries\color{white} Description} \\
\midrule
Name & Display name. \\
Description & Free-text notes. \\
Coordinates (JSON) & The GeoJSON coordinate array for the polygon. Polygons require an outer ring (first and last coordinate identical to close the ring). \\
Geometry (JSON) & Advanced: full GeoJSON geometry object. \\
Properties (JSON) & Advanced: additional GeoJSON feature properties. \\
\bottomrule
\end{longtable}

\begin{notebox}
SAR search areas created by the POB feature appear in the Regions tab automatically.
Their \code{properties} include the \code{pobUuid} linking them to the originating POB incident.
\end{notebox}

% ── Chapter 16: Notes ─────────────────────────────────────────
\chapter{Notes}
\label{ch:notes}

Notes are free-text annotations attached to geographic positions.
They function as sticky notes on the chart --- useful for recording pilotage tips, local hazards not
on the chart, marina details, or crew observations.

\screenshot{mydata-notes-list}{My Data --- Notes tab, list view with title, description, href, MIME type, timestamp}{}

\section{Notes Table Columns}

Notes have a distinct column set reflecting their attachment-oriented nature:

\rowcolors{2}{LtGray}{white}
\begin{longtable}{L{2.4cm} L{11.3cm}}
\toprule
\rowcolor{Navy}
{\sffamily\bfseries\color{white} Column} &
{\sffamily\bfseries\color{white} Description} \\
\midrule
Title & The display name of the note (the first column is labelled \textit{Title} for notes, not \textit{Name}, to reflect the note-taking context). \\
Description & Summary text. For imported text files, this is set to the first 120 characters of the file content. \\
Href & An optional file path or URL linking the note to an attachment --- a PDF pilot book page, an image, a text file. When present, tapping \key{Open} opens the attachment directly in the built-in viewer. \\
MIME Type & The MIME type of the attachment, e.g.\ \code{application/pdf}, \code{image/jpeg}, \code{text/plain}. Used to route the file to the correct viewer. \\
Timestamp & ISO 8601 creation or last-update timestamp. \\
\bottomrule
\end{longtable}

\section{Creating a Note}

Tap \key{New}. The Note Editor drawer opens.


\rowcolors{2}{LtGray}{white}
\begin{longtable}{L{3cm} L{10.7cm}}
\toprule
\rowcolor{Navy}
{\sffamily\bfseries\color{white} Field} &
{\sffamily\bfseries\color{white} Description} \\
\midrule
Title & The note title (displayed in the \textit{Title} column). \\
Description & Main text content of the note. \\
Href & Optional path or URL to an attachment file. Enter an absolute path (e.g.\ \path{/home/pi/charts/harbour.pdf}) or a URL. \\
MIME Type & Optional MIME type for the attachment. FairWindSK uses this to select the correct viewer. If left blank, FairWindSK detects the type from the file extension. \\
Has position & Check this box to attach a geographic position to the note. \\
Latitude & Latitude when \textit{Has position} is checked. \\
Longitude & Longitude when \textit{Has position} is checked. \\
\bottomrule
\end{longtable}

\section{Opening a Note}

\subsection{Notes Without an Attachment}

Tapping \key{Open} opens the note editor in read-only mode, showing the title, description, and position.

\subsection{Notes With an Attachment (href set)}

Tapping \key{Open} immediately opens the attachment:

\begin{itemize}
  \item \textbf{Local file that exists:} opens in the built-in FileViewer (PDF, image, or text viewer depending on the MIME type).
  \item \textbf{HTTP/HTTPS URL:} opens in an embedded web view inside the Bottom Bar horizontal drawer.
  \item \textbf{Local directory path:} shows a message directing you to use My Data~> Files to browse the folder (single-window model prevents opening a separate file manager).
\end{itemize}

\section{Importing Notes}

The Notes tab accepts a broader set of import file types than other tabs.
In addition to GeoJSON/JSON:

\begin{itemize}
  \item \textbf{Text files} (\code{.txt}, \code{.md}, \code{.html}): FairWindSK creates a note with the filename as the title and the first 120 characters of the file content as the description, setting the \code{href} to the absolute file path and detecting the MIME type automatically.
  \item \textbf{Images} (\code{.jpg}, \code{.jpeg}, \code{.png}): Creates a note with the filename as title and sets the \code{href} and \code{mimeType} accordingly.
  \item \textbf{PDF files}: Creates a note pointing to the PDF file via \code{href}.
\end{itemize}

The import filter text reads: \textit{Notes and attachments (*.geojson *.json *.txt *.md *.markdown *.html *.htm *.pdf *.jpg *.jpeg *.png)}.

% ── Chapter 17: Charts ────────────────────────────────────────
\chapter{Charts}
\label{ch:charts}

The Charts tab lists available chart tile sources registered on the Signal~K server.
It manages the chart source catalogue --- not the actual chart display, which happens inside
Freeboard-SK or another chart plotter running in the Application Area.

\screenshot{mydata-charts-list}{My Data --- Charts tab, list view with name, format, identifier, URL, timestamp}{}

\section{Charts Table Columns}

\rowcolors{2}{LtGray}{white}
\begin{longtable}{L{2.4cm} L{11.3cm}}
\toprule
\rowcolor{Navy}
{\sffamily\bfseries\color{white} Column} &
{\sffamily\bfseries\color{white} Description} \\
\midrule
Name & Display name of the chart source. \\
Description & Free-text description. \\
Format & The chart tile format. Supported formats include: \code{mbtiles} (MBTiles SQLite database), \code{kap} (BSB/KAP raster charts), \code{gemf} (GEMF tile archive), \code{tilemapresource} (TileMapResource XML), \code{package} (ZIP chart package). \\
Identifier & The chart identifier used by Signal~K and chart plotters to reference this source. \\
Chart URL & The local filesystem path or URL where the chart tiles are served. \\
Timestamp & ISO 8601 registration timestamp. \\
\bottomrule
\end{longtable}

\section{Adding a Chart Source}

Tap \key{New}. The Chart Editor drawer opens.

\rowcolors{2}{LtGray}{white}
\begin{longtable}{L{3cm} L{10.7cm}}
\toprule
\rowcolor{Navy}
{\sffamily\bfseries\color{white} Field} &
{\sffamily\bfseries\color{white} Description} \\
\midrule
Name & Display name for this chart source. \\
Description & Free-text description. \\
Identifier & Unique identifier string used by chart plotters to reference this source. \\
Chart format & The format of the chart tiles. FairWindSK detects the format from the file extension when importing. \\
Chart URL & Path or URL to the chart tiles. For MBTiles, this is the absolute path to the \code{.mbtiles} file. \\
Tilemap URL & For \code{tilemapresource} format only: the path to the \code{tilemapresource.xml} file. The Chart URL is then set to the directory containing it. \\
Chart region & Optional: a GeoJSON bounding box or geometry describing the geographic coverage of this chart. \\
Chart scale & Optional: the chart scale denominator (e.g.\ \code{50000} for 1:50,000). \\
Chart layers & Optional: comma-separated list of layer names for tile sources that support multiple layers. \\
Chart bounds & Optional: GeoJSON coordinates describing the chart bounds. \\
\bottomrule
\end{longtable}

\section{Importing a Chart Source}

Tap \key{Import}. Select a chart file.
FairWindSK detects the chart format from the file extension and constructs a chart resource object
automatically.
For \code{tilemapresource.xml} files, FairWindSK sets both the tile URL (the directory) and the
tilemap URL (the XML file path).

The import creates or updates the chart record on the Signal~K server.

\section{Exporting the Chart Catalogue}

Tap \key{Export}. The chart catalogue is exported as a JSON manifest document with structure:

\begin{lstlisting}
{
  "type": "FairWindSKChartPackage",
  "charts": [ { ... }, { ... } ]
}
\end{lstlisting}

This manifest can be imported on another FairWindSK installation to replicate the chart catalogue
configuration (it describes chart sources, not the chart tile data itself).

% ── Chapter 18: Tracks ────────────────────────────────────────
\chapter{Tracks --- History}
\label{ch:tracks}

The Tracks tab shows the historical GPS track of the vessel --- the recorded sequence of position
fixes over time.
Tracks are generated automatically by Signal~K when the tracks plugin is enabled.

\screenshot{mydata-tracks-list}{My Data --- Tracks tab, table of track points with time, lat/lon, altitude}{}

\section{Track Table Columns}

The track table shows individual position fixes (points) in chronological order:

\rowcolors{2}{LtGray}{white}
\begin{longtable}{L{2.4cm} L{11.3cm}}
\toprule
\rowcolor{Navy}
{\sffamily\bfseries\color{white} Column} &
{\sffamily\bfseries\color{white} Description} \\
\midrule
Time & The timestamp of the position fix, formatted in local time using the system locale. \\
Latitude & Latitude to 6 decimal places. \\
Longitude & Longitude to 6 decimal places. \\
Altitude & Altitude in metres (if available from the GPS fix; otherwise shown as 0). \\
\bottomrule
\end{longtable}

\section{Live Track Recording}

If Signal~K is actively providing position updates, new track points are appended to the table in
real time.
The table scrolls to show the most recent point automatically.

\section{Viewing a Track Point}

Select a row and tap \key{Open}. A read-only dialog opens showing the timestamp, latitude, longitude,
and altitude for that point in full detail.

\subsection{Editing a Track Point}

Select a row and tap \key{Edit}. The Track Point Editor opens.

\rowcolors{2}{LtGray}{white}
\begin{longtable}{L{3cm} L{10.7cm}}
\toprule
\rowcolor{Navy}
{\sffamily\bfseries\color{white} Field} &
{\sffamily\bfseries\color{white} Description} \\
\midrule
Timestamp & Date and time of the fix. Uses a date/time picker with up/down arrows. \\
Latitude & Latitude (8 decimal places, range $-90$ to $+90$). \\
Longitude & Longitude (8 decimal places, range $-180$ to $+180$). \\
Altitude & Altitude in metres. \\
\bottomrule
\end{longtable}

\section{Adding a Track Point Manually}

Tap \key{New}. The same editor opens with default values (current time and position 0,0).
Enter the correct values and tap \key{OK}.
The point is inserted into the track in chronological order.

\section{Deleting a Track Point}

Select a row and tap \key{Delete}. Confirm when prompted. The point is removed from the track.

\section{Exporting the Track}

Tap \key{Export} to save the complete track as a GeoJSON \code{LineString} or
\code{FeatureCollection}:

\begin{lstlisting}
{
  "type": "FeatureCollection",
  "features": [
    {
      "type": "Feature",
      "geometry": {
        "type": "LineString",
        "coordinates": [[lon, lat, alt], ...]
      }
    }
  ]
}
\end{lstlisting}

The default export filename is \path{tracks-history.geojson}.

\begin{tipbox}
Export the track after every passage and store it with the voyage logbook.
The GeoJSON file can be imported into most GIS tools, chart plotters, and voyage-planning software
for review and analysis.
\end{tipbox}

\section{Importing a Track}

Tap \key{Import} and select a GeoJSON file containing a \code{LineString} geometry.
FairWindSK extracts each coordinate in the line as a track point with a synthetic timestamp,
appending them to the current track in order.

\section{Uses of the Track Tab}

\begin{itemize}
  \item \textbf{Passage review:} compare the actual track against the planned route. Identify deviations, traffic avoidance manoeuvres, and current effects.
  \item \textbf{Anchorage assessment:} see the swing circle after anchoring overnight.
  \item \textbf{Safety investigation:} reconstruct the vessel's path before and during an incident for formal reporting.
  \item \textbf{Logbook cross-reference:} correlate positions at specific times with logbook entries.
\end{itemize}

% ── Chapter 19: Files ─────────────────────────────────────────
\chapter{Files}
\label{ch:files}

The Files tab is a direct filesystem browser for files and directories accessible on the device
running FairWindSK.
It provides navigation, search, viewing, and file management operations (copy, cut, paste, rename,
delete, new folder) without leaving the FairWindSK interface.

\screenshot{mydata-files-browser}{My Data --- Files tab, directory browser with toolbar}{}

\section{Interface Layout}

The Files tab is divided into three areas:

\begin{itemize}
  \item \textbf{Toolbar} at the top: navigation and action buttons, path bar, and search field.
  \item \textbf{File list} in the centre: the current directory contents as a tree view with columns for Name, Size, and Date. Row height is 58\,px for comfortable touch selection. Alternating row colours aid readability.
  \item \textbf{Actions row} below the toolbar: file operation buttons that appear when items are selected.
\end{itemize}

\section{Navigation}

\subsection{Directory Navigation}

\rowcolors{2}{LtGray}{white}
\begin{longtable}{L{3cm} L{10.7cm}}
\toprule
\rowcolor{Navy}
{\sffamily\bfseries\color{white} Action} &
{\sffamily\bfseries\color{white} How to trigger} \\
\midrule
Enter a folder & Double-tap a folder row in the file list, or select and tap \key{Open}. \\
Go up one level & Tap \key{↑} (Up) button in the toolbar. \\
Navigate by path & Tap the path bar at the top of the toolbar, type an absolute path, and press Return. \\
Go home & Tap \key{Home} button. Returns to the home directory of the running user. \\
\bottomrule
\end{longtable}

\subsection{Column Sorting}

Tap any column header to sort by that column.
The sort direction toggles between ascending and descending on subsequent taps.
The default sort is by Name, ascending.

\section{Searching Files}

\subsection{Basic Search}

Tap \key{Search} or press Return in the search field.
FairWindSK recursively searches the current directory tree for files and subdirectories whose names
contain the search text.

Results are shown in a flat list in the file view.
Tap a result to select it; double-tap to open.
The search continues to the root of the browsed directory tree.


\subsection{Search Options}

Tap \key{Options} (or the filter button) to expand the search options:

\rowcolors{2}{LtGray}{white}
\begin{longtable}{L{3.5cm} L{10.2cm}}
\toprule
\rowcolor{Navy}
{\sffamily\bfseries\color{white} Option} &
{\sffamily\bfseries\color{white} Description} \\
\midrule
Case Sensitive & When enabled, the search matches the exact capitalisation of the search text. Default: off (case-insensitive). \\
Search Hidden & When enabled, hidden files and directories (names beginning with \code{.}) are included in search results. Default: off. \\
Search System & When enabled, system files and directories (e.g.\ \code{/proc}, \code{/sys} on Linux) are included. Default: off. \\
\bottomrule
\end{longtable}

The search options button label updates to show a compact summary of active options (e.g.\
\textit{Options: case, hidden}).

\section{File Actions}

\subsection{Opening Files}

Select a file and tap \key{Open}, or double-tap the row.

FairWindSK detects the file type using MIME type detection and routes it to the appropriate viewer:

\rowcolors{2}{LtGray}{white}
\begin{longtable}{L{4cm} L{9.7cm}}
\toprule
\rowcolor{Navy}
{\sffamily\bfseries\color{white} File type} &
{\sffamily\bfseries\color{white} Viewer} \\
\midrule
Images (\code{.jpg}, \code{.jpeg}, \code{.png}, \code{.gif}) & Built-in image viewer with pinch-to-zoom and pan. \\
PDF (\code{.pdf}) & Built-in PDF viewer, page by page. \\
Text files (\code{.txt}, \code{.md}, \code{.log}, \code{.html}) & Built-in text viewer, formatted as plain text or HTML. \\
Other file types & A message is shown; the file cannot be previewed but can be managed. \\
\bottomrule
\end{longtable}

Directories open in the directory browser (navigating into them) rather than in a viewer.
Attempting to open a directory path from a Notes href shows a message directing you to use the Files
tab to browse it instead (single-window model constraint).


\subsection{Renaming Files and Folders}

Select an item and tap \key{Rename}.
An input dialog opens in the Bottom Bar horizontal drawer.
Enter the new name and tap \key{OK}.
The file or directory is renamed on the filesystem immediately.

\subsection{Creating a New Folder}

Tap \key{New Folder}.
An input dialog opens.
Enter the folder name and tap \key{OK}.
The new folder is created in the current directory.

\subsection{Copy, Cut, and Paste}

\begin{enumerate}
  \item Select one or more files or folders.
  \item Tap \key{Copy} to mark them for copying, or \key{Cut} to mark them for moving.
  \item Navigate to the destination directory.
  \item Tap \key{Paste}.
\end{enumerate}

FairWindSK performs safety checks before paste operations:

\begin{itemize}
  \item Copying a directory into itself or one of its subdirectories is blocked.
  \item Pasting into the source directory of a cut operation shows a warning.
  \item The self-paste safety check runs a self-test at startup to verify the detection logic is working correctly.
\end{itemize}

Directory copy and move operations are recursive --- the entire subtree is copied or moved.

\subsection{Deleting Files and Folders}

Select one or more items and tap \key{Delete}.
A confirmation dialog opens in the Bottom Bar horizontal drawer.
Confirm to delete.

\begin{warningbox}
File deletion is permanent.
There is no recycle bin or undo for file operations.
Take particular care when deleting directories --- all contents are deleted recursively.
\end{warningbox}

\section{File List Columns}

The file list shows three columns:

\rowcolors{2}{LtGray}{white}
\begin{longtable}{L{2cm} L{11.7cm}}
\toprule
\rowcolor{Navy}
{\sffamily\bfseries\color{white} Column} &
{\sffamily\bfseries\color{white} Description} \\
\midrule
Name & Filename or directory name, with an icon indicating the file type (folder, image, document, etc.). \\
Size & File size in human-readable units (B, KB, MB, GB). Directories show the size of their immediate contents. \\
Date & Last modification timestamp. The \textit{Type} column (column index 2) is hidden by default. \\
\bottomrule
\end{longtable}

\section{Practical Uses of the Files Tab}

\begin{itemize}
  \item Browse pilot book excerpts stored as PDFs on the Raspberry Pi.
  \item View passage planning images (harbour photos, aerial views, approach charts).
  \item Access voyage log files or NMEA recordings.
  \item Manage chart tile archives before registering them in the Charts tab.
  \item Review diagnostic log files generated by FairWindSK (see Chapter~\ref{ch:settings-system}).
\end{itemize}
